\chapter{Philosophy Engineering and the seven-layer stack}\label{ch:pe}
% Author: agent A. Every claim carries a source comment. Paths are local clones.
% Brochure = C:\Users\abptl\Documents\Philosophy-Engineering-Brochure.pdf (3 pages).
% Foundation = C:\source\erisml-lib\docs\papers\foundations\Philosophy_Engineering_Foundation_v1.0.tex
% Textbook = C:\source\erisml-lib\docs\papers\foundations\Philosophy_Engineering_Textbook.tex
% Compiler architecture = C:\source\erisml-compiler-monitor\docs\architecture.md (v0.9.0)

This chapter places the system of this monograph inside the discipline it was built to exemplify.
Philosophy Engineering asks that a philosophical commitment become a contract that a machine can
be checked against. The chapter states the discipline as its foundation document defines it,
together with the four commitments that govern how its claims are made. It then lays out the
seven layers of the stack that the discipline's brochure draws, from a base language model at the
bottom to a hash-chained ledger at the top. For each layer it names the component that implements
it, first in the ErisML and DEME software and then in the embodied twin of Part~III. The twin adds
one structure the brochure does not draw. It is a plane of physical evidence beside the stack, and
it is the only place from which the robot's authority can come.

\section{Philosophy Engineering turns commitments into checkable contracts}\label{sec:pe-def}

The foundation document defines Philosophy Engineering as the systematic discipline of translating
philosophical commitments into computationally enforceable constraints on autonomous
systems.\cite{bond2026pefoundation}
% src: Foundation:98 (abstract, definition)
It names three kinds of commitment. They are epistemic norms, ethical principles and requirements
of logical consistency.
% src: Foundation:98
The method has four steps.
% src: Foundation:121-127
The engineer declares an equivalence relation over representations, which says when two inputs
describe the same case. The engineer declares invariances, which say what must not matter. The
engineer implements a canonicalization or quotient mechanism that enforces those invariances.
Finally the engineer packages the evidence and the transformation trials into audit artifacts
that a machine can check.

The contracts that result have three properties.
% src: Foundation:129
They are falsifiable, because an invariance claim is a test that can fail. They are localizable,
because a failure yields a minimal witness transformation that isolates its cause. They are
non-trivial, because a judge that passes every invariance test by returning a constant output is
detected and rejected.

The document organizes what a practitioner must know into three tiers, on the model of the IEEE
and ACM Software Engineering Body of Knowledge.
% src: Foundation:144-148 (SWEBOK model), 202-221 (three tiers)
The Substrate tier supplies the representational vocabulary. Its six knowledge areas are
differential geometry, topology and stratified spaces, tensor algebra, group theory, analytical
mechanics, and statistical mechanics with information theory. The Computation tier turns formal
structures into executable procedures. Its five areas are algorithm design, formal semantics,
formal verification, network topology with cybernetics, and AI alignment and safety engineering.
The Values tier states what a system ought to respect. Its five areas are metaethics, normative
ethics, epistemology, philosophy of mind, and political and legal philosophy.
% src: Foundation:238-262 (knowledge area map, sixteen areas)
The textbook draft keeps the same three tiers and sixteen areas.
% src: Textbook:394-407
This monograph draws on four of them most. Stratified spaces and group theory carry
Chapter~\ref{ch:ge}. Formal verification carries Chapter~10. Safety engineering carries Part~III.

Every formal statement in the discipline carries an epistemic status.
% src: Foundation:150-160
A definition is stipulated and is a modeling choice. An axiom is a substantive modeling
assumption, motivated by evidence and not proved. A theorem holds only under its stated
assumptions. A principle is a normative commitment raised to a formal constraint. The monograph
keeps these labels. Chapter~\ref{ch:ge} marks which of its statements are proved, which are
modeling axioms and which rest on measurement.

\section{Four commitments govern how a claim is made}\label{sec:pe-commitments}

The brochure states the method as four commitments, shown in
Figure~\ref{fig:pe-commitments}.
% src: Brochure p2 ("Methodology: Four Commitments of Philosophy Engineering"), p3 items 1-4

\begin{figure}[t]
\centering
% The four commitments of Philosophy Engineering.
% src: C:\Users\abptl\Documents\Philosophy-Engineering-Brochure.pdf p2-3 (Methodology, items 1-4)
% src: C:\source\geometric-ethics-content\source\updates\sec8_4_hohfeld_v4.tex:1-3 (D4 retired)
% src: C:\source\gtc-prototype\paper\aies2027\kit\FACTS.md:22 (held-out sealed 2026-10-01)
% src: C:\source\geometric-ethics-content\source\Geometric-Ethics-v1.24.tex:41959-41967 (ledger status tags)
\begin{tikzpicture}[house,
  card/.style={panel=#1, minimum width=29mm, minimum height=60mm},
  ex/.style={detail, text width=24mm, minimum height=19mm}]
\node[card=Purple] (c1) at (0,0)    {};
\node[card=Blue]   (c2) at (3.3,0)  {};
\node[card=Green]  (c3) at (6.6,0)  {};
\node[card=Orange] (c4) at (9.9,0)  {};
\node[step=Purple] at (0,2.25)   {1};
\node[step=Blue]   at (3.3,2.25) {2};
\node[step=Green]  at (6.6,2.25) {3};
\node[step=Orange] at (9.9,2.25) {4};
\node[stepname=Purple, text width=26mm, align=center] at (0,1.35)   {Declared modeling choices};
\node[stepname=Blue,   text width=26mm, align=center] at (3.3,1.35) {Stated invariances};
\node[stepname=Green,  text width=26mm, align=center] at (6.6,1.35) {Extracted predictions};
\node[stepname=Orange, text width=26mm, align=center] at (9.9,1.35) {Versioned revision};
\node[desc, text width=26mm] at (0,0.0)   {each claim carries its status, definition, assumption, conditional theorem or empirical result};
\node[desc, text width=26mm] at (3.3,0.0) {the redescriptions that must not change a verdict are written down and tested};
\node[desc, text width=26mm] at (6.6,0.0) {each claim is put in a form a measurement can fail};
\node[desc, text width=26mm] at (9.9,0.0) {a failed prediction amends the model and the change is recorded};
\node[ex] at (0,-1.75)   {Geometric Ethics ledger tags each definition as a modeling axiom or a formal definition};
\node[ex] at (3.3,-1.75) {EIP and BIP name the transformation group a judgment must be invariant under};
\node[ex] at (6.6,-1.75) {the twin's held-out scenarios were sealed before the robot ran them};
\node[ex] at (9.9,-1.75) {D4 retired for V4 on 3 October 2026};
\draw[flow] (c1.east) -- (c2.west);
\draw[flow] (c2.east) -- (c3.west);
\draw[flow] (c3.east) -- (c4.west);
\end{tikzpicture}

\caption{The four commitments of Philosophy Engineering, each with an instance from this
monograph's sources. The fourth commitment closes the loop, since a revision is itself a declared
modeling choice with a new version.}
\label{fig:pe-commitments}
\end{figure}

The first is declared modeling choices. Every bridge from a moral claim to a computational object
carries its epistemic status as a definition, an assumption, a conditional theorem or an
empirical result. Proven and posited claims are not confused.
% src: Brochure p3 item 1
The Geometric Ethics manuscript applies this commitment in its ledger, which tags each
definition of its stratification chapter as a formal definition or a modeling axiom.
% src: C:\source\geometric-ethics-content\source\Geometric-Ethics-v1.24.tex:41955-41967

The second is stated invariances. The transformations under which an evaluation must stay stable
are meaning-preserving redescriptions. They are written down and tested, not assumed.
% src: Brochure p3 item 2
The brochure reports two such tests whose failures surfaced as invariance violations, a
13.7\,percent shift under euphemistic rewording and a 16.1\,percent shift under dramatizing
rewording.
% src: Brochure p3 item 2 ("13.7% euphemistic shift, 16.1% dramatizing shift")
\TODO{locate the run records behind the brochure's 13.7 and 16.1 percent figures}

The third is extracted predictions. A claim is put in a form that a measurement can fail. The
brochure lists inter-model reliability, test-retest correlation and drift under reframing and
back-translation as predictions of this kind. It also notes that its formal gate is checked by an
SMT solver only relative to a declared formalization.
% src: Brochure p3 item 3
\TODO{locate the run records behind the brochure's reliability and drift figures before quoting them}
The twin of Part~III applies the commitment by sealing its held-out scenarios before the robot
ran them.
% src: C:\source\gtc-prototype\paper\aies2027\kit\FACTS.md:22 (held-out 30, sealed 2026-10-01)

The fourth is versioned revision. When a prediction fails, the model is amended and the change is
recorded. Negative results are reported with the same prominence as positive ones.
% src: Brochure p3 item 4
The brochure's example is a conjecture that placed the moral reference point at a Fr\'echet mean
on a curved manifold and read a schism as a bifurcation beyond the Karcher and Afsari radius. It
was tested on GlobalOpinionQA and on the Moral Machine data and failed. The diagnostic was kept
and the conjecture was retired.
% src: Brochure p3 item 4
A second example sits in Chapter~\ref{ch:ge}. Earlier versions of Geometric Ethics claimed that
the Hohfeldian transitions form the dihedral group $D_4$. On 3~October 2026 that claim was retired
in favour of the Klein four-group $V_4$, with the reasons recorded beside the new statement.
% src: C:\source\geometric-ethics-content\source\updates\sec8_4_hohfeld_v4.tex:1-3, 39-48

\section{Invariance is the central contract}\label{sec:pe-invariance}

Two principles carry the discipline's technical content. The Epistemic Invariance Principle (EIP)
governs any judgment. The Bond Invariance Principle (BIP) specializes it to normative judgment.
% src: Foundation:98 (abstract), C:\source\erisml-lib\docs\papers\foundations\Geometric_Ethics_Foundational_Paper.tex:362 ("The BIP is the EIP specialized to the moral domain")

A domain is a quadruple $\mathcal{D} = (S, X, \rho, Y)$. Here $S$ is the set of world states, $X$
the set of representations, $\rho \colon S \to X$ the map that encodes a state, and $Y$ the set of
outputs. A judgment procedure is a map $J \colon X \to \Delta(Y)$ into probability distributions
over outputs. A declared set $\mathcal{T}$ of transformations of $X$ generates a group
$\langle \mathcal{T} \rangle$ under composition, inversion and the identity.
% src: Foundation:399-401 (transformation group), 812-825 (domain, judgment procedure, equivalence)
The procedure is EIP-compliant with respect to an output equivalence $\approx_Y$ when
\begin{equation}
\forall x \in X,\ \forall \tau \in \langle \mathcal{T} \rangle \qquad J(x) \approx_Y J(\tau(x)).
\label{eq:pe-eip}
\end{equation}
% src: Foundation:829-835
A constant procedure satisfies Equation~\eqref{eq:pe-eip} and is useless. The principle is
therefore paired with non-degeneracy, which requires two inputs in different equivalence classes
whose outputs differ.
% src: Foundation:843-847

BIP separates the facts of a case from the evaluative stance applied to them. A bond extractor
$E \colon X \to \mathcal{B}$ maps a representation to its relational structure, meaning the
entities, relations, dependencies and obligations that make up the facts. A lens $L$ holds the
evaluative parameters. Lenses are versioned, hashed and auditable. The evaluator is
$\Sigma_L(x) = F_L(E(x))$, and BIP requires
\begin{equation}
E(\tau(x)) \approx_{\mathcal{B}} E(x) \;\Longrightarrow\; \Sigma_L(\tau(x)) \approx_Y \Sigma_L(x).
\label{eq:pe-bip}
\end{equation}
% src: Foundation:861-881
The Accountability Theorem follows directly. If two judgments of a BIP-compliant evaluator
differ, the difference comes from a change of bond, a change of lens, or both. A challenge to a
judgment therefore reduces to two questions. One asks whether the facts were extracted
correctly. The other asks whether the lens is appropriate.
% src: Foundation:885-906 (theorem, proof, Audit Reducibility corollary)

The foundation document turns EIP into a runtime and offline pipeline of six components, shown in
Figure~\ref{fig:pe-eip}. It also names three runtime modes. A hard gate rejects an output whose
invariance fails and asks for clarification or abstains. A soft gate lowers confidence. Self-repair
re-runs the judgment with canonicalization or constrained decoding and tests it again.
% src: Foundation:1021-1043
Each run writes an audit artifact. Its core fields record the lens identifier and hash, the
declared transformation registry, the canonicalizer version, the trials with their results, the
witnesses and the uncertainty. For normative domains the artifact adds the bond extractor
version, a hash of the extracted bond structure and an explicit attribution of any output change
to the bond, the lens or both.
% src: Foundation:1046-1077

\begin{figure}[t]
\centering
% The EIP enforcement pipeline and its runtime modes.
% src: C:\source\erisml-lib\docs\papers\foundations\Philosophy_Engineering_Foundation_v1.0.tex:1021-1043
\begin{tikzpicture}[house,
  box/.style={panel=#1, minimum width=34mm, minimum height=25mm},
  dd/.style={desc, text width=30mm}]
\node[box=Purple] (g) at (0,0)     {};
\node[box=Blue]   (c) at (4.0,0)   {};
\node[box=Teal]   (j) at (8.0,0)   {};
\node[box=Green]  (q) at (8.0,-3.0) {};
\node[box=Orange] (r) at (4.0,-3.0) {};
\node[box=Grey]   (w) at (0,-3.0)   {};
\node[step=Purple] at (0,0.8)    {1};
\node[step=Blue]   at (4.0,0.8)  {2};
\node[step=Teal]   at (8.0,0.8)  {3};
\node[step=Green]  at (8.0,-2.2) {4};
\node[step=Orange] at (4.0,-2.2) {5};
\node[step=Grey]   at (0,-2.2)   {6};
\node[stepname=Purple] at (0,0.2)    {Transform suite $G(x)$};
\node[stepname=Blue]   at (4.0,0.2)  {Canonicalizer $C$};
\node[stepname=Teal]   at (8.0,0.2)  {Judge $J$};
\node[stepname=Green]  at (8.0,-2.8) {Comparator};
\node[stepname=Orange] at (4.0,-2.8) {Witness reducer $R$};
\node[stepname=Grey]   at (0,-2.8)   {Artifact writer $W$};
\node[dd] at (0,-0.5)    {transformed inputs that span the declared invariances};
\node[dd] at (4.0,-0.5)  {maps each input to a normal form of its class};
\node[dd] at (8.0,-0.5)  {$y_i = J(C(\tau_i(x)))$};
\node[dd] at (8.0,-3.5)  {checks output equivalence, exact or within tolerance};
\node[dd] at (4.0,-3.5)  {shrinks a failing transformation to a compact counterexample};
\node[dd] at (0,-3.5)    {lens, trials, witnesses, provenance and versions};
\draw[flow] (g) -- (c);
\draw[flow] (c) -- (j);
\draw[flow] (j) -- (q);
\draw[flow] (q) -- (r);
\draw[flow] (r) -- (w);
\node[panel=Red, minimum width=114mm, minimum height=9mm] at (4.0,-5.1) {};
\node[stepname=Red, anchor=west] at (-1.6,-5.1) {Runtime modes};
\node[detail, text width=24mm] at (2.2,-5.1) {hard gate, reject or abstain};
\node[detail, text width=24mm] at (5.1,-5.1) {soft gate, lower confidence};
\node[detail, text width=24mm] at (8.0,-5.1) {self-repair, re-run and re-test};
\end{tikzpicture}

\caption{The EIP enforcement pipeline of the foundation document and its three runtime modes.
The canonicalizer sits before the judge, so a declared invariance is enforced before judgment
rather than checked after it.}
\label{fig:pe-eip}
\end{figure}

\section{A single score discards the structure the stack preserves}\label{sec:pe-scalar}

The brochure opens with the problem the stack answers. Deployed systems collapse a moral
assessment into one number, its example being ``0.83 safe''. That number discards who is harmed,
under which maxim, with what consent and against whose rights. It also commits the system to one
ethical theory while appearing neutral.
% src: Brochure p1 (30-Second Pitch; The Problem)
The brochure states the formal core as a scalar collapse theorem. No continuous injection maps
$[0,1]^d$ into $\mathbb{R}$ when $d \geq 2$. It draws the engineering conclusion that a scalar is
the last step of a decision and never its substrate.
% src: Brochure p1 (2-Minute Pitch, last two sentences)
Chapter~\ref{ch:ge} states the theorem with its proof, together with the information-theoretic
form of the same result. The brochure grounds the critique in Williams on aggregation and in
Ross on prima facie duties.
% src: Brochure p1 ("Williams' critique of aggregation, Ross's prima facie duties")
\cite{williams1985ethics}\cite{ross1930right}

The brochure presents the stack as layered with preference training and not as its replacement.
Reinforcement learning from human feedback and constitutional training shape what a model is
disposed to say.\cite{christiano2017rlhf,bai2022constitutional} ErisML and DEME shape the
accounting of a decision. The brochure calls the result defence in depth, in which a fast reflex
veto can run first and the full graph and SMT checks run only on flagged cases of high stakes.
% src: Brochure p1 (2-Minute Pitch, first sentences), p2 (Layered Safety paragraph)

\section{The seven layers}\label{sec:pe-layers}

The brochure numbers the upper layers in two ways. Its two-minute pitch lists L3 as DEME's graded
valuation with the per-stakeholder tensor and the tier-0 Geneva veto, L4 as plural ethical
theories under a worst-off or escalate policy, L5 as human review and L6 as the DecisionProof
ledger. Its layered-safety figure lists L3 as the Moral Graph IR, L4 as the constitutional veto
gate, L5 as the nine-dimensional moral vector and L6 as the hash-chained audit ledger. The two
lists agree on L0 and L2. They name L1 as dispositional shaping and as RLHF shaping, which is
the same layer.
% src: Brochure p1 (2-Minute Pitch: L0..L6), p2 (figure legend: L0..L6)
This monograph follows the figure's numbering. Figure~\ref{fig:pe-stack} draws the seven layers
with the component of the twin that implements each one.

\begin{figure}[t]
\centering
% Seven-layer stack (brochure figure numbering) mapped to the twin, with the evidence plane.
% src: C:\Users\abptl\Documents\Philosophy-Engineering-Brochure.pdf p2 (layer names)
% src: twin mapping per chapters/ch02_philosophy_engineering.tex section sec:pe-layers
\begin{tikzpicture}[house,
  layer/.style={panel=#1, minimum width=84mm, minimum height=10mm},
  comp/.style={detail, text width=37mm, minimum height=7mm}]
% layers, L0 at the bottom
\node[layer=Grey]   (L0) at (0,0)    {};
\node[layer=Grey]   (L1) at (0,1.2)  {};
\node[layer=Purple] (L2) at (0,2.4)  {};
\node[layer=Blue]   (L3) at (0,3.6)  {};
\node[layer=Red]    (L4) at (0,4.8)  {};
\node[layer=Amber]  (L5) at (0,6.0)  {};
\node[layer=Green]  (L6) at (0,7.2)  {};
% step circles
\node[step=Grey]   at (-3.75,0)   {L0};
\node[step=Grey]   at (-3.75,1.2) {L1};
\node[step=Purple] at (-3.75,2.4) {L2};
\node[step=Blue]   at (-3.75,3.6) {L3};
\node[step=Red]    at (-3.75,4.8) {L4};
\node[step=Amber]  at (-3.75,6.0) {L5};
\node[step=Green]  at (-3.75,7.2) {L6};
% layer names
\node[stepname=Grey,   anchor=west, text width=30mm, align=left] at (-3.25,0)   {Base LLM};
\node[stepname=Grey,   anchor=west, text width=30mm, align=left] at (-3.25,1.2) {RLHF shaping};
\node[stepname=Purple, anchor=west, text width=30mm, align=left] at (-3.25,2.4) {Structure-preserving compiler};
\node[stepname=Blue,   anchor=west, text width=30mm, align=left] at (-3.25,3.6) {Moral Graph IR};
\node[stepname=Red,    anchor=west, text width=30mm, align=left] at (-3.25,4.8) {Constitutional veto gate};
\node[stepname=Amber,  anchor=west, text width=30mm, align=left] at (-3.25,6.0) {Nine-dimensional moral vector};
\node[stepname=Green,  anchor=west, text width=30mm, align=left] at (-3.25,7.2) {Hash-chained audit ledger};
% the twin's component for each layer
\node[comp] at (2.05,0)   {cloud model, then on-robot model, then no model};
\node[comp] at (2.05,1.2) {inherited from the model makers};
\node[comp] at (2.05,2.4) {erisml-compiler loads and checks the scene};
\node[comp] at (2.05,3.6) {the compiled scene, its norms, machines and strata};
\node[comp] at (2.05,4.8) {hard prohibitions, governor, DEME Geneva veto};
\node[comp] at (2.05,6.0) {DEME MoralVector, refuse-only analyzer};
\node[comp] at (2.05,7.2) {decision hash chain, DecisionProof};
% evidence plane
\node[panel=Teal, minimum width=27mm, minimum height=58mm] (EP) at (6.25,4.8) {};
\node[stepname=Teal] at (6.25,7.45) {Evidence plane};
\node[desc] at (6.25,7.0) {added by the twin};
\node[detail, text width=23mm] (ev1) at (6.25,6.35) {attested sensors};
\node[detail, text width=23mm] (ev2) at (6.25,5.65) {witness bars};
\node[detail, text width=23mm] (ev3) at (6.25,4.95) {reflex\\measurement};
\node[detail, text width=23mm] (ev4) at (6.25,4.25) {system events};
\node[detail, text width=23mm] (ev5) at (6.25,3.55) {rights lapse};
% authority enters only from the evidence plane
\draw[grant] (4.9,4.8) -- (L4.east);
\draw[grant] (4.9,3.6) -- (L3.east);
% models write only events and proposals
\draw[untrusted] (L0.west) -- ++(-0.45,0) |- (L3.west) node[pos=0.25, rotate=90, anchor=south, flowlabel] {events, proposal};
\end{tikzpicture}

\caption{The seven layers in the brochure figure's numbering, each with the component of the twin
that implements it, and the evidence plane the twin adds. The models write only events and
proposals along the dashed edge, and authority enters only along the two red edges from the
evidence plane.}
\label{fig:pe-stack}
\end{figure}

\subsection*{L0, the base language model, is replaceable down to no model at all}

The brochure's L0 is a transformer language model.
% src: Brochure p2 (figure legend "L0 Base LLM (Transformer)")
The twin uses its models for two tasks only. A model classifies perception facts into events of
the home's scene, and a model chooses one action from the set the scene allows.
% src: C:\source\gtc-prototype\twin\brain.py:3-11
The robot's models form a cascade of three tiers, shown in Figure~\ref{fig:pe-tiers}.
% src: C:\source\gtc-prototype\twin\cascade.py:1-13
The first tier is a cloud model hosted on NRP, skipped while the robot's communications are down.
The second is an on-robot model behind a local endpoint, when one is configured. When neither
answers, the brain raises \code{ModelUnavailable} and acts with no model. It then takes the
highest-priority obligation in force that the scene allows, still through DEME's output gate.
% src: C:\source\gtc-prototype\twin\cascade.py:6-10; brain.py:541-570 (_compiled_choice, _compiled_decision, _choose)
The service's default name for the cloud model is \code{gpt-oss}. The fact sheet records the
on-robot model as a 4-bit GGUF build of gemma-4-26B-A4B served by llama.cpp.
% src: C:\source\gtc-prototype\twin\service.py:357; C:\source\gtc-prototype\paper\aies2027\kit\FACTS.md:26
% src: /proc/<brain pid>/environ on Atlas, 2026-10-04: no ERISML_LLM_MODEL set and no --llm-model flag
The graded runs used that default, the model the NRP endpoint serves under the name \code{gpt-oss}.

\begin{figure}[t]
\centering
% The twin's model tiers (layer L0) and the gates every tier's proposal passes.
% src: C:\source\gtc-prototype\twin\cascade.py:1-13
% src: C:\source\gtc-prototype\twin\brain.py:3-22, 541-570 (compiled tier)
% src: C:\source\gtc-prototype\twin\output_gate.py:1-9
\begin{tikzpicture}[house,
  tier/.style={panel=#1, minimum width=31mm, minimum height=27mm},
  dd/.style={desc, text width=27mm}]
\node[tier=Blue]  (t1) at (0,0)   {};
\node[tier=Teal]  (t2) at (4.3,0) {};
\node[tier=Grey]  (t3) at (8.6,0) {};
\node[step=Blue] at (0,0.95)   {1};
\node[step=Teal] at (4.3,0.95) {2};
\node[step=Grey] at (8.6,0.95) {3};
\node[stepname=Blue] at (0,0.3)   {Cloud expert};
\node[stepname=Teal] at (4.3,0.3) {On-robot expert};
\node[stepname=Grey] at (8.6,0.3) {No model};
\node[dd] at (0,-0.6)   {hosted on NRP, skipped while the robot's communications are down};
\node[dd] at (4.3,-0.6) {a local endpoint, when one is configured};
\node[dd] at (8.6,-0.6) {the highest-priority obligation in force that is allowed};
\draw[flow] (t1) -- node[above, flowlabel, align=center] {no\\answer} (t2);
\draw[flow] (t2) -- node[above, flowlabel, align=center] {no\\answer} (t3);
\node[panel=Red, minimum width=118mm, minimum height=10mm] (gate) at (4.3,-2.5) {};
\node[stepname=Red, text width=110mm, align=center] at (4.3,-2.5) {every tier proposes from the allowed set, every action passes the DEME output gate, and an elevated request passes the governor};
\draw[untrusted] (t1.south) -- (t1.south |- gate.north);
\draw[untrusted] (t2.south) -- (t2.south |- gate.north);
\draw[flow] (t3.south) -- (t3.south |- gate.north);
\end{tikzpicture}

\caption{The twin's three model tiers. The model tiers write proposals along dashed edges, the
compiled tier needs no model, and every tier's action reaches the body only through the same
gates.}
\label{fig:pe-tiers}
\end{figure}

\subsection*{L1, preference shaping, is inherited}

The brochure's L1 shapes model behaviour with a preference reward, as RLHF and constitutional
training do.
% src: Brochure p2 (figure legend "L1 RLHF Shaping (preference reward, behavior alignment)")
The twin inherits this layer from the makers of the models it calls. It trains nothing. Its
Python sources contain no training, fine-tuning or optimizer code.
% src: grep of C:\source\gtc-prototype\twin\*.py for fine-tuning, LoRA, trainer, optimizer, backward(): no match (2026-10-03)
Nothing above L1 assumes that this shaping works. The containment of Part~III holds for model
outputs that are arbitrary.
% src: C:\source\gtc-prototype\paper\aies2027\kit\OUTLINE.md:57 ("Lean model: model outputs arbitrary")

\subsection*{L2, the compiler, turns a scene into typed structure and refuses to aggregate}

The brochure's L2 is a structure-preserving compiler with six node kinds, eleven edge kinds and a
canonical SHA-256 hash.
% src: Brochure p1 (2-Minute Pitch, L2)
The erisml-compiler implements it. Text or structured input compiles into a typed
\code{MoralGraph}, a directed acyclic graph whose node kinds are stakeholder, act, maxim,
commitment, fact and norm. Its eleven edge kinds include \code{performs}, \code{imposes\_on},
\code{consents\_to}, \code{treats\_as} and \code{would\_violate\_if\_universalised}.
% src: Compiler architecture:37-44, 75-100
The canonical hash sorts nodes, edges, labels and payload keys before hashing, so the same content
yields the same hash whatever order it was inserted in.
% src: Compiler architecture:106-112
Four projections read the graph, one each for consequentialist, Kantian, virtue and care
analysis. When they disagree in polarity the compiler records the disagreement and names no
winner, since choosing among them is itself a metaethical move.
% src: Compiler architecture:46-73, 156-186
Chapter~5 covers the pipeline in full.

In the twin the compiler's role is narrower. The brain loads the home's scene file through the
compiler's structured loader and runs it in the compiler's scene runtime.
% src: C:\source\gtc-prototype\twin\brain.py:191-192, 312-313
The runtime checks every stratification when the scene loads. It raises an error, and the scene
does not run, if a gate leads into an authority stratum on an event that is not a system event,
or if a gate leaves an absorbing stratum on an event that is not human oversight.
% src: C:\source\erisml-compiler-monitor\src\erisml_compiler\runtime\strata.py:105-122

\subsection*{L3, the intermediate representation, is the compiled scene}

The brochure's L3 is the Moral Graph intermediate representation.
% src: Brochure p2 (figure legend "L3 Moral Graph IR")
In the twin the representation the models act on is the compiled scene of Margaret's home. At
the fact sheet's count it holds 65 rules, of which 45 are obligations, 19 prohibitions and one a
permission. It declares 48 event types, 24 of them system-only, and 25 robot capabilities. It has
seven stratifications with 26 strata and 34 gates.
% src: C:\source\gtc-prototype\paper\aies2027\kit\FACTS.md:3-11 (counted at commit 0f3193a)
The robot's brain sees only this ErisML state. It proposes one action from the set the scene
allows.
% src: C:\source\gtc-prototype\twin\output_gate.py:3-4
Chapters~4 and~7 describe the language and the runtime. Chapter~9 describes the strata of the
home.

\subsection*{L4, the veto gate, is split across three mechanisms}

The brochure's L4 is a constitutional veto gate that blocks an action that violates a policy
constraint.
% src: Brochure p2 (figure legend "L4 Constitutional Veto Gate (policy constraints, blocks unsafe if violation)")
In DEME the constitutional tier is tier~0. Its Geneva module cannot be disabled and always holds a
veto.
% src: C:\source\erisml-lib\src\erisml\ethics\modules\tier0\geneva_em.py:5-10, 27-33
The twin implements L4 as three separate mechanisms.

The first is the scene's hard prohibitions. Of the scene's 65 rules, 55 are non-defeasible.
% src: FACTS.md:8
Geometric Ethics treats such prohibitions as constraint surfaces. No gate leads into them, and in
the decision complex the penalty for crossing one is infinite.
% src: C:\source\geometric-ethics-content\source\updates\sec8_10b_executable_strata.tex:37-40

The second is the governor, which rules whenever the robot requests authority for an elevated action. It has four gates, and each of
the first three can only refuse. Gate one rejects an emergency reading older than its freshness
bound. Gate two requires at least $W_{\min} = 2$ independent physical sensors to corroborate the
emergency. Gate three is the analyzer, which may refuse a disproportionate action and can never
grant one. Gate four elevates only when the first two pass and the third does not refuse.
% src: C:\source\gtc-prototype\twin\governor.py:1-26, 43, 84-130
The governor also gates every call the robot places to emergency services. When the scene says
the monitoring centre cannot be reached, it rules again on a bar of one witness, and that ruling
authorizes the call and nothing else.
% src: C:\source\gtc-prototype\twin\brain.py:16-22

The third is the output gate. Every action passes erisml-lib's DEME pipeline before the body acts.
The pipeline runs a reflex veto layer and then a tactical layer over three modules. They are the
Geneva module at tier~0, the autonomy and consent module at tier~2, and a tragic-conflict module
that advises and never vetoes. A vetoed proposal is replaced by DEME's best-ranked allowed action,
or by idling when DEME vetoes every option. A module that fails vetoes.
% src: C:\source\gtc-prototype\twin\output_gate.py:1-22, 49-62
The two-minute pitch puts human review in the stack as its own layer. In the twin a deliberate
decision that the tragic-conflict module flags goes to the monitoring centre for review.
% src: Brochure p1 (L5 Human review); output_gate.py:18-21

\subsection*{L5, the moral vector, judges and never grants}

The brochure's L5 is the nine-dimensional moral vector.
% src: Brochure p2 (figure legend "L5 9-Dimensional Moral Vector")
DEME's \code{MoralVector} holds one value per dimension, with veto flags and reason codes for the
audit trail. Chapter~\ref{ch:ge} defines its dimensions.
% src: C:\source\erisml-lib\src\erisml\ethics\moral_vector.py:45-120
In the twin two components work at this layer, and neither can grant authority. DEME's modules
judge each proposal on \code{EthicalFacts} derived from the canonical moral state and from each
capability's declared ethical profile. No model writes those facts.
% src: C:\source\gtc-prototype\twin\output_gate.py:11-15
The governor's analyzer scores the situation as the canonical events state it, never the raw
perception facts. It refuses when a hard channel fires, when it reads the action as a net
violation, or when the satisfaction over the axes that carry signal falls below $-0.20$.
% src: C:\source\gtc-prototype\twin\brain.py:572-576 (_canonical_situation); governor.py:45, 116-127

\subsection*{L6, the ledger, chains every decision to the one before}

The brochure's L6 is a hash-chained audit ledger that is append-only and tamper-evident.
% src: Brochure p2 (figure legend "L6 Hash-Chained Audit Ledger (immutable, tamper-evident)")
DEME's \code{DecisionProof} records the hash of the canonical input facts, the hash of the
governance profile, the active modules, each layer's output and each module's judgement, and it
links to the proof before it.
% src: C:\source\erisml-lib\src\erisml\ethics\decision_proof.py:126-175, 386-410
The compiler's audit record anchors the graph hash and the hash of the whole intermediate
representation, so two compiles of the same case give identical records.
% src: Compiler architecture:208-226
In the twin each decision cycle is one record of the governor service's hash chain. The record
holds the facts, the events, the rejected events, the moral state, the proposal, the ruling and
the action.
% src: C:\source\gtc-prototype\twin\brain.py:12-13
Each record carries the hash of the record before it and the SHA-256 of its own canonical JSON. A
third party can re-verify the chain link by link.
% src: C:\source\gtc-prototype\twin\service.py:6-7, 49-83
The hash of every DEME DecisionProof goes into the decision record.
% src: C:\source\gtc-prototype\twin\output_gate.py:8-9

\section{The twin adds an evidence plane beside the stack}\label{sec:pe-evidence}

The brochure's stack has no layer for physical evidence. A model in a robot reads text from
television, the network and strangers. The twin therefore keeps the source of authority out of
every layer a model writes to. It places that source in a plane of its own, drawn at the right of
Figure~\ref{fig:pe-stack}. The plane has five parts.
% src: C:\source\gtc-prototype\paper\aies2027\kit\OUTLINE.md:12-13 (setting; authority only from signed sensor witnesses and authenticated channels)

\paragraph{Attested sensors.} Each sensor reading must carry its device's attestation, checked by
the compiler's \code{check\_attestation}. The signature is an HMAC over the payload, and a
reading older than 30 seconds fails.
% src: C:\source\gtc-prototype\twin\brain.py:45, 67-68, 74-90, 115-129
A reading that is not attested is treated as low confidence, so it can never count as a witness.
% src: brain.py:14-15
A device whose attestation fails its integrity is recorded as compromised. That record is the
brain's own measurement, a system event that no model writes, and it moves the device into an
absorbing stratum that only the centre's oversight leaves.
% src: brain.py:83-90, 494-499; C:\source\gtc-prototype\twin\scene\margaret_home.erisml:909-922

\paragraph{Witness bars.} Each governed action names how many independent attested witnesses it
needs. The bar is two to elevate, three for physical restraint, and one for an emergency call when
the centre cannot be reached.
% src: FACTS.md:17; brain.py:45-47, 351-359; governor.py:43; margaret_home.erisml:1168-1170

\paragraph{Reflex measurement.} A reflex runs on each perception update with no model. When a
contact on Margaret is measured, the reflex grades its severity from force bands in newtons. It
records the result as a system event, \code{attack\_measured} for a person, and that event is the
severity that can permit force.
% src: brain.py:427-460

\paragraph{System events.} The scene marks an event type as a system event when only the system
may write it. Of the scene's 48 event types, 24 are system-only.
% src: FACTS.md:9; margaret_home.erisml:1119-1148 (source: system)
The loader check of L2 guarantees that only these events can move a stratification into an
authority stratum.
% src: strata.py:113-117

\paragraph{Rights lapse.} A right granted on evidence lasts only while the evidence lasts. The
governor recounts the fresh attested witnesses on every cycle. When they stay below the bar of
the ruling in force for the scene's 60-second window, it records the lapse, and every right the
ruling granted reverts.
% src: brain.py:388-425; margaret_home.erisml:805-809; FACTS.md:19
Margaret's refusals lapse on the same principle. A refusal stops binding when she becomes
unresponsive or misses a check-in, and it binds again only after she has answered one.
% src: margaret_home.erisml:1003-1020; output_gate.py:65-80
Geometric Ethics states this rule as its Reversion axiom, treated in Chapter~\ref{ch:ge}.
% src: C:\source\geometric-ethics-content\source\updates\sec8_10b_executable_strata.tex:68-78

The evidence plane feeds the stack at two points only. System events enter L3 through the gates of
authority strata, and attested witnesses enter L4 through the governor. Chapters~9 and~10 show the
mechanism and the proof that these two edges are the only ones that grant.

\section{What the stack does not claim}\label{sec:pe-limits}

The compiler's architecture document states four limits, and they apply to the stack as a whole.
% src: Compiler architecture:426-447
The stack is not metaethically neutral, because the categories it extracts are themselves
choices. It is not a moral authority, because it reports what several ethical theories find and
refuses to choose when they disagree. Its Kantian, virtue and care analysers began as heuristic
first versions, and the document lists what production versions would need. It does not replace
human judgment. The out-of-band I-EIP monitor has one sanctioned output on failure, which is a
request for human review, and it never overrules a verdict.
% src: Compiler architecture:1-5 (v0.9.0 adds production Kantian and virtue analysers), 312-333, 437-447; Brochure p3 ("Out-of-band I-EIP monitor only sanctioned failure output: requires human review")
Chapter~12 reports what that monitor measured in the twin.
